\newpage
\section{模型的建立与求解}

\subsection{发酵动力学模型}
啤酒发酵过程遵循质量作用定律，建立如下反应动力学模型：
$$\frac{dC_S}{dt} = -k_1 C_S - k_2 C_S C_X$$
$$\frac{dC_E}{dt} = Y_{E/S} k_1 C_S$$
$$\frac{dC_{Ac}}{dt} = k_3 C_S - k_4 C_{Ac}$$
其中$C_S$为糖度，$C_E$为乙醇浓度，$C_{Ac}$为乙酸浓度，$C_X$为菌体浓度，$k_1, k_2, k_3, k_4$为反应速率常数。

\subsection{参数估计}
采用Levenberg-Marquardt算法拟合实验数据，目标函数为：
$$\min_{\theta} \sum_{t} \|C^{\text{model}}(t; \theta) - C^{\text{exp}}(t)\|^2$$
迭代终止条件：$\|\theta^{(k+1)} - \theta^{(k)}\| < 10^{-6}$。

\subsection{模型求解结果}
估计得到的速率常数为：$k_1=0.023$, $k_2=0.001$, $k_3=0.005$, $k_4=0.002$（单位：min$^{-1}$或L/(mol·min)）。交叉验证RMSE=0.034，模型拟合良好。
